\section{Composition and the Entity Subspace}
\label{sec:compose}

The composition task arranges $K$ entities in one cycle $\pi$.
The entities are represented by orthonormal vectors $e_1,\ldots,e_K$.
The one-step binding
is $k_i=e_i$, $v_i=e_{\pi(i)}$, so the output uses the same vector
space as the next input. A query at depth $h$ asks for $e_{\pi^h(i)}$,
obtained by applying the same trained operator $h$ times: $Z^h e_i$.
There are no learned
per-hop parameters. Training uses depths $\{1, 2, 3\}$; evaluation
adds $\{4, 5, 6\}$ and probes $\{7, 21\}$ ($d = 16$, $K = 8$,
orthonormal keys). A single full cycle avoids the shorter cycles of
a general random permutation, which can turn nominally held-out
depths into identity or training queries. Depth 21 has the same
target as depth 5 in an 8-cycle; it tests stability under additional
applications (Appendix~\ref{app:period}).

\textbf{Recovery across depths.} Four of five unconstrained seeds
reach $\recninety \geq 0.9996$ at every tested depth; three reach
1.000 throughout. At the training budget, the fifth seed reaches
0.93 at $h{=}1$, falling to 0.16 at $h{=}7$ and 0.0001 at $h{=}21$.
% <!-- evidence: R3 -->

\begin{figure}[htbp]
\centering
\includegraphics[width=\textwidth]{fig_depth.pdf}
\caption{Recovery under repeated application of the learned matrix
($d{=}16$, $K{=}8$). \emph{Left:} mean cosine for the converged seed
with rank cap $k{=}7{=}K{-}1$, compared with a prediction computed
from its entity-subspace operator and the ideal cycle. \emph{Right:} the fraction of queries with
cosine greater than 0.9, averaged over the four converged
unconstrained seeds (near 1.0), and shown separately for the
rank-capped seed and the fifth unconstrained seed.}
% <!-- evidence: R4 -->
% <!-- evidence: R3 -->
\label{fig:depth}
\end{figure}

\textbf{Error accumulation.} With rank cap $k = K - 1$, exact linear
recovery is impossible. Nevertheless, the one converged seed reaches
mean cosine 0.92--0.93 at shallow depths ($h \leq 5$).
Its $\recninety$ falls from 0.996 at $h = 1$ to 0.881 at $h = 7$
and 0.060 at $h = 21$ (Figure~\ref{fig:depth}).
% <!-- evidence: R4 -->
A calculation using the entity-subspace operator and ideal cycle
predicts measured cosine within 0.008 through $h = 7$ (Table~\ref{tab:depthcurve}).
% <!-- evidence: R4 -->

\textbf{Action on the entity subspace.} The converged solutions have
whole-matrix effective rank 16.0 at $d{=}16$ and stable rank
14.6--15.6. These whole-matrix measures include components outside
the subspace used by the task.
% <!-- evidence: R5 -->
We therefore decompose $Z$ into blocks relative to the ideal cycle's
$K$-dimensional entity subspace $\mathcal{E}$. The block acting within
$\mathcal{E}$ has effective rank 7.999--8.000 (maximum 8) and a
scale-corrected residual of 0.7--2.4\% from the ideal cycle. Leakage
out of $\mathcal{E}$ is below 1.5\% of the state norm. The complement
block is full-rank and non-contractive
(Table~\ref{tab:subspace}, Appendix~\ref{app:subspace}).
% <!-- evidence: R5 -->
These measurements support an approximately invariant entity
subspace on which effective rank is close to $K$. Rank-capped seeds
that fail recovery also fail the corresponding subspace test.
